\chapter{Introduction}
\label{cha:introduction}
As software becomes more critical,
and the amount of programming languages increases \missingcitation \todo{Provide a source or statistic for the proliferation of programming languages}, the need to ensure
an instruction in a program does what the programmer think it does becomes essential \missingcitation \todo{Cite a source discussing the 'software crisis' or the importance of software reliability}. \elias{En lite krånglig första mening}
This is especially important when compilers and interpreters translate code into
lower level representations such as machine code or bytecode \hampus{true, but} \missingcitation \todo{Cite a source on compiler correctness or the risks of translation errors, e.g., CompCert}. \elias{Eftersom du inte gör det kanske det inte är något du ska trycka på precis i början}

To bridge this gap of syntax and semantics, the field of programming language \todo{is the term "gap" applicable here} \elias{Lite oklart vad som avses här}
semantics was developed, formalizing the behavior of instructions
in a rigorous mathematical representation \missingcitation \todo{Cite a foundational text on formal semantics, e.g., Plotkin, Kahn, or Winskel}. \hampus{was it developed for this reason?} \elias{Jag tror att du ska prata om det som något som finns, snarare än något som har uppfunnits. Kanske ha vinkeln att alla programspråk har någon semantik och att programmerare har en informell uppfattning om den, men att formell semantik möjliggör formella bevis på korrekthet av program och analyser}.
By defining these meanings, we can reason about a program's behavior and
validate that it works as intended regardless of input \missingcitation \todo{Cite the benefits of formal methods in verification}.

In academic settings, such as the course
"Semantics of Programming Languages (1DL311)" provided by Uppsala University,
these definitions and proofs are often presented "on paper"~\citep{uppsala_syllabus_2023, uppsala_reading_2023}.
These paper-based proofs frequently omit inductive steps that share
similarities with previous ones \hampus{maybe give a example / page number in the book}.
While these proofs are useful for
teaching, it is also susceptible to human error, implicit assumptions and
gaps in logical reasoning \missingcitation \todo{Cite research on the fallibility of hand-written proofs vs. machine-checked proofs, e.g., De Millo et al. or Wiedijk}. By mechanizing \elias{Inte definierat!} the proof we introduce a layer of
automatic validation, thereby ensuring logical soundness and verify
formal correctness \hampus{Maybe cite Lean / Coq / general proof assistant for this statement} \missingcitation \todo{Cite the general utility of Proof Assistants/LCF-style kernels}. \elias{…men man kan inte heller säga ”liknar föregående fall” eller ”inses lätt”!}.

This thesis investigates the process of translating these pen-and-paper proofs
into mechanized proofs using the proof assistant Lean~\citep{theorem_proving_2026}. The goal is to
evaluate whether the techniques used in Semantics literature~\citep{uppsala_reading_2023} correspond
directly to Lean's mechanisms or if they require any reformulation to achieve
machine-level precision.

%
%
%
\label{intro:research_question}
\section{Research Question and Goals}
The primary objective of this work is to determine the "semantic gap"
between human-readable proofs and machine-checked logic.
Specifically, this thesis seeks to answer the following research questions:
\begin{itemize}
    \item \textbf{RQ1:} Do the proof techniques presented in the 1DL311 curriculum
    map directly to Lean's internal mechanisms, or do they require significant
    structural reformulation?
    \item \textbf{RQ2:} What implicit assumptions or "pedagogical shortcuts"
    present in the literature must be made explicit to satisfy the Lean kernel?
\end{itemize}

%
%
%
\label{intro:motivation}
\section{Motivation}
\elias{Remove subsection and combine with 1.0} 
Informal proofs are designed for human consumption and often prioritize
intuition over exhaustive detail \hampus{this statement feels weired especially "humna consumption"}.
However, these omissions can hide underlying complexities or logical leaps \hampus{no proof of this} \missingcitation \todo{Cite a source discussing 'social processes' in proofs or the 'De Bruijn factor'}.
By identifying these points of divergence,
this research provides a roadmap for how educators can
structure pen-and-paper proofs to decrease ambiguity and better prepare students
for formal verification tasks.

%
%
%
\section{Scope}
\label{intro:scope}
This thesis will evaluate proofs from the course "Semantics of Programming Languages
(1DL311)" which uses the literature "Semantics with Applications: An Appetizer"~\citep{nielson_semantics_2007}.
will be limited to the evaluation to a few proofs from the chapter
"1. Introduction" ~\citep[p.~1]{nielson_semantics_2007} and
"2. Operational Semantics" ~\citep[p.~19]{nielson_semantics_2007}.
These chapters introduce the toy programming language "While", which will be
used in the proofs mechanized into Lean 4.

%
%
%
\section{Contribution}
\label{intro:contribution}
This thesis contribution to the field of programming language is the process of
mechanizing written programming language semantic proofs into Lean 4,
as well as which assumptions if any is needed to be stated explicitly in the proof
assistant that otherwise was implicitly defined in the written counterpart.

%
%
%
\label{intro:roadmap}
\section{Thesis Structure Roadmap} 
\todo{A map of the thesis layout written in plain English to give reader a tour guide through the thesis} \noindent
\RED{\lipsum[1]}